\documentclass[11pt]{article}
\usepackage{explorer}
\newcommand{\ExplorerLanguage}{FR}
\newcommand{\ExplorerNumber}{02}
\newcommand{\ExplorerKind}{Leçon et méthode}
\begin{document}
\ExplorerTitle{Seconde · olympiades}{Le principe des tiroirs}
\ExplorerSection{Avant de commencer}
Division euclidienne et dénombrement élémentaire.\par
\ExplorerSection{Idées et méthode}
Si \(n + 1\) objets sont rangés dans n tiroirs, au moins un tiroir contient deux objets. Sinon, chaque tiroir en contiendrait au plus un et il y aurait au plus n objets.\par
Plus généralement, si N objets occupent k tiroirs, un tiroir en contient au moins ⌈\(\frac{N}{k}\)⌉. Le plafond ⌈x⌉ est le plus petit entier supérieur ou égal à x.\par
Le choix des tiroirs fait la preuve : restes d’une division, mois de naissance, petits carrés d’un quadrillage ou classes de nombres. Il faut justifier que chaque objet appartient à un tiroir.\par
\ExplorerSection{Exemple expliqué}
Parmi 13 personnes, deux sont nées le même mois : les 12 mois sont les tiroirs. Cela ne garantit pas deux anniversaires le même jour.\par
\ExplorerSection{À explorer}
Choisir six entiers distincts entre 1 et 10. Les grouper par paires de somme 11 et repérer une paire complète.\par
\end{document}
